\section{Server Load and Curved Spacetime (Complexity)}

\textit{(Server Load and Curved Spacetime - Complexity)}

\begin{quote}
\textbf{``Matter tells spacetime how to curve? No, data load tells the server how to allocate computing power. When you enter a capital city, the screen lags, character models load slowly---that's not spacetime curvature; that's the server's computational load being too high in that area. Gravity is the `damping' produced when computational systems process complex information.''}
\end{quote}

In previous chapters, we established: light speed is bandwidth, space is projection. Now, we face physics' biggest BOSS---\textbf{Gravity}.

In Einstein's general relativity, gravity is explained as spacetime curvature. But he didn't explain \textbf{why it curves}.

In \textbf{Code of Azeroth} theory, we will demystify gravity. We will demonstrate: Gravity is not a force that pulls objects together, but an \textbf{entropic force}. Simply put, it's a \textbf{lag effect} produced when the server processes complex data.

\subsection{Gravity Is Emergent: Like Air Pressure}

To understand gravity, think about air pressure. If you look at individual molecules, there's no concept of ``pressure.'' Only when you count countless molecular collisions do you get pressure.

Similarly, gravity is the statistical manifestation of spacetime's microscopic bits (quantum pixels).

\textbf{Definition 5.1.1 (Entropic Force)}

Matter tends to move toward places of low gravitational potential (high curvature) because this movement maximizes the system's \textbf{disorder (entropy)}. Or, this is the computational system finding the \textbf{minimum computational cost path}.

\subsection{Complexity Equals Volume: Graphics Card Killer}

Frontier research on holographic principles (CV conjecture) tells us: \textbf{Spatial volume equals computational complexity.}

\[
V \sim \text{Complexity}
\]

\begin{itemize}
\item A region looking ``large'' means the system needs to execute many lines of code to render this region.
\item The volume inside a black hole grows with time, meaning the black hole's data structure becomes increasingly chaotic; decompressing it requires increasing computing power over time.
\end{itemize}

Therefore, \textbf{curved spacetime} is actually a \textbf{``server load heatmap''}.

\subsection{Clock Slowing Due to Excessive Load}

Now we can answer: Why do massive objects (stars/black holes) warp spacetime?

\begin{enumerate}
\item \textbf{Mass as complexity}: A massive object is essentially a highly entangled, data-heavy \textbf{high-complexity texture}. For the server, this is a ``graphics card killer.''
\item \textbf{Computational black hole}: To render this complex object, the server must allocate massive CPU cycles to this region.
\item \textbf{Processing delay (gravitational redshift)}: According to \textbf{computational conservation law}, if the server is busy calculating this object's internal structure, its speed of processing external messages (like photons passing by) slows down.
\end{enumerate}

\textbf{Conclusion}: From the player's perspective, ``time'' in that region slows down. When light passes there, due to the processing node's \textbf{network congestion}, the path deflects (gravitational lensing).

\subsection{Gravity Is Route Redirection}

Objects ``fall'' toward massive objects because on that path, \textbf{information transmission delay is minimized}.

Imagine you're playing an online game.
\begin{itemize}
\item Flat spacetime = empty wilderness, walking is straight.
\item Curved spacetime = crowded capital city.
\end{itemize}

When you approach the capital, due to data packet congestion, the system automatically adjusts your route, directing you toward that load center (though this sounds counterintuitive, in four-dimensional spacetime's geodesic equations, this is the shortest path).

\textbf{Gravity is actually route redirection caused by network congestion.}

\subsection{Geometric Deformation in Tensor Networks}

We can demonstrate gravity using \textbf{tensor networks}.

Imagine a flat fishing net (flat space).
Now we want to place a super complex object on this net. To encode this object's data, we need to insert more \textbf{nodes} into the original grid.

According to the CV conjecture, inserting more nodes equals increasing that region's ``volume.'' But with fixed boundaries, increasing internal volume forces the grid to \textbf{bulge (curve)}.

\textbf{Conclusion}:

Einstein's field equations are actually the holographic computer's \textbf{resource scheduling equations}:
\begin{itemize}
\item Left side $G_{\mu\nu}$ (curvature): Represents \textbf{distribution of computational nodes}.
\item Right side $T_{\mu\nu}$ (energy): Represents \textbf{data load to be processed}.
\end{itemize}

The equation shows: To process high-density data load ($T_{\mu\nu}$), the system must reconstruct network topology ($G_{\mu\nu}$), increasing node density, thus causing macroscopic spacetime curvature.

\textbf{Gravity is the `noise' emitted when the universe computer runs at full load.}
